\editionPage{defi-cor-1-1}{Solution du défi D01}{Deux laboratoires}{Une démarche possible ; d’autres rédactions sont recevables.}
\begin{questions}
\item $u_A=0,08/4=\qty{0,02}{mm}$ ; $u_B=0,12/6=\qty{0,02}{mm}$.
\item B a des mesures individuelles plus dispersées, mais ses mesures plus nombreuses conduisent ici à la même incertitude sur la moyenne.
\item Les écarts de 0,12 et 0,10 mm dépassent nettement les incertitudes annoncées. Rechercher une origine commune : étalonnage ou protocole. La répétition ne suffit pas à la corriger.
\end{questions}
\editionRetenir{La qualité d’une mesure dépend à la fois de sa dispersion et de la maîtrise des décalages systématiques.}

\editionPage{defi-cor-2-1}{Solution du défi D02}{Contrôler une étiquette}{Une démarche possible ; d’autres rédactions sont recevables.}
\begin{questions}
\item $m=270,0\unit{g}$ ; $\rho=270,0/250=\qty{1,08}{g.mL^{-1}}$.
\item $100\times40,5/270,0=15,0\,\%$ : l’annonce massique est respectée.
\item La concentration en masse vaut $40,5/0,250=\qty{162}{g.L^{-1}}$, donc 16,2 g pour 100 mL. Masse de solution et volume de solution ne se confondent pas.
\end{questions}
\editionRetenir{Un pourcentage massique et une concentration en masse répondent à deux questions différentes.}

\editionPage{defi-cor-3-1}{Solution du défi D03}{Préparer avec le stock disponible}{Une démarche possible ; d’autres rédactions sont recevables.}
\begin{questions}
\item $V_p=6,00\times250,0/30,0=\qty{50,0}{mL}$.
\item Effectuer deux prélèvements de 25,0 mL, les transférer dans la même fiole, puis compléter au trait et homogénéiser. Chaque prélèvement doit suivre le mode d’emploi de la pipette.
\item Non : un seul prélèvement donnerait $30,0\times25,0/250,0=\qty{3,00}{g.L^{-1}}$.
\end{questions}
\editionRetenir{La dilution conserve le soluté prélevé ; le volume final est celui de la solution entière.}

\editionPage{defi-cor-4-1}{Solution du défi D04}{Reconstituer une entité}{Une démarche possible ; d’autres rédactions sont recevables.}
\begin{questions}
\item Il possède $16-8=8$ protons : c’est l’oxygène, $\ce{^{16}_{8}O}$.
\item La charge négative traduit deux électrons supplémentaires : 10 électrons.
\item Deux ions sodium par ion oxyde : $\ce{Na2O}$.
\end{questions}
\editionRetenir{Le noyau identifie l’élément ; le cortège électronique détermine la charge de l’ion.}

\editionPage{defi-cor-5-1}{Solution du défi D05}{Retrouver un élément}{Une démarche possible ; d’autres rédactions sont recevables.}
\begin{questions}
\item $1s^2\,2s^2\,2p^6\,3s^2\,3p^5$ : 17 électrons, c’est le chlore.
\item Il gagne un électron : $\ce{Cl-}$, avec la configuration de l’argon.
\item Le chlorure possède 17 protons et une charge négative ; l’argon en possède 18 et est neutre. Leurs noyaux diffèrent.
\end{questions}
\editionRetenir{Une configuration de valence stable explique la formation de l’ion ; elle ne change pas l’identité du noyau.}

\editionPage{defi-cor-6-1}{Solution du défi D06}{Une solution vue de deux échelles}{Une démarche possible ; d’autres rédactions sont recevables.}
\begin{questions}
\item $m_e=6m_C+12m_H+6m_O=\qty{2,9904e-25}{kg}$.
\item $m=12,0\times0,250=\qty{3,00}{g}$ ; $N=3,00\times10^{-3}/m_e\simeq\num{1,00e22}$.
\item $n=N/N_A\simeq\qty{0,0167}{mol}$. On a compté uniquement le glucose, sans inclure l’eau.
\end{questions}
\editionRetenir{Les quotients $m/m_e$ et $N/N_A$ relient la masse macroscopique au comptage microscopique.}

\editionPage{defi-cor-7-1}{Solution du défi D07}{Une réserve de froid}{Une démarche possible ; d’autres rédactions sont recevables.}
\begin{questions}
\item $E=0,300\times3,34\times10^5=\qty{1,002e5}{J}$.
\item $t=100200/25=\qty{4008}{s}$, soit environ 67 min, sous l’hypothèse d’un apport constant.
\item Après disparition de la glace, l’eau liquide se réchauffe. Le calcul de fusion seul ne détermine pas cette élévation de température.
\end{questions}
\editionRetenir{Un changement d’état peut échanger de l’énergie sans variation de température du corps pur.}

\editionPage{defi-cor-8-1}{Solution du défi D08}{Choisir le bon mélange}{Une démarche possible ; d’autres rédactions sont recevables.}
\begin{questions}
\item Il faut 0,300 mol de \ce{H2} et 0,150 mol de \ce{O2}.
\item $0,300/2<0,300/1$ : \ce{H2} limite. Il se forme 0,300 mol d’eau et il reste 0,150 mol de \ce{O2}.
\item La quantité d’eau est identique ; le deuxième mélange gaspille un excès de dioxygène.
\end{questions}
\editionRetenir{Une quantité de produit visée ne suffit pas à garantir des proportions stœchiométriques ; il faut régler les deux réactifs.}

\editionPage{defi-cor-9-1}{Solution du défi D09}{Le produit est-il prêt ?}{Une démarche possible ; d’autres rédactions sont recevables.}
\schema[corrige]{ccm-synthese}
\begin{questions}
\item La tache commune avec N soutient cette hypothèse ; compléter l’identification si nécessaire.
\item Non : une autre espèce est détectée, compatible avec du réactif restant.
\item Après une purification adaptée au protocole, refaire une CCM avec S, N et R sur la même plaque. La disparition de la tache de R appuie l’amélioration de la pureté.
\end{questions}
\editionRetenir{Identifier le produit recherché et vérifier sa pureté sont deux étapes distinctes.}

\editionPage{defi-cor-10-1}{Solution du défi D10}{Auditer une réaction}{Une démarche possible ; d’autres rédactions sont recevables.}
\begin{questions}
\item $x=236-140-94=2$ ; $92=54+38$.
\item Une fission : un noyau lourd donne deux noyaux de tailles plus petites.
\item $N=9,0\times10^7/(3,0\times10^{-11})=\num{3,0e18}$ événements.
\end{questions}
\editionRetenir{Une équation nucléaire s’analyse par les noyaux et les conservations, pas par une combustion ordinaire.}

\editionPage{defi-cor-11-1}{Solution du défi D11}{Valider un pointage}{Une démarche possible ; d’autres rédactions sont recevables.}
\schema[corrige]{chronophoto}
\begin{questions}
\item Les valeurs sont 1 ; 3 ; 5 ; 7 m/s, toutes orientées vers la droite.
\item Les longueurs sont 1,5 cm et 3,5 cm ; la figure corrigée montre les vecteurs au bon point d’application.
\item Les différences des valeurs calculées sont régulières, mais les images ne donnent qu’un échantillonnage et des approximations. La conclusion doit mentionner le modèle et les incertitudes.
\end{questions}
\editionRetenir{Un pointage fournit un modèle discret du mouvement, avec une durée et une échelle spatiale à vérifier.}

\editionPage{defi-cor-12-1}{Solution du défi D12}{Une pesanteur plus faible}{Une démarche possible ; d’autres rédactions sont recevables.}
\begin{questions}
\item La distance est doublée ; l’attraction est divisée par quatre.
\item $P=80\times9,8/4=\qty{196}{N}$, soit environ 200 N.
\item Non : l’attraction gravitationnelle agit sur les objets en orbite. Une sensation d’apesanteur n’est pas une disparition de cette interaction.
\end{questions}
\editionRetenir{Le poids dépend du lieu ; la gravitation ne disparaît pas dès que l’on quitte le sol.}

\editionPage{defi-cor-13-1}{Solution du défi D13}{Trois phases d’un lancer}{Une démarche possible ; d’autres rédactions sont recevables.}
\schema[corrige]{chute}
\begin{questions}
\item Poids vers le bas partout ; vitesse vers le haut en montée, nulle au sommet, vers le bas en descente.
\item Le poids n’est pas compensé et la vitesse change de part et d’autre du sommet.
\item Étalonner les positions et le temps, tracer les vitesses, vérifier la trajectoire et rechercher des indices de frottement. Le modèle s’évalue sur plusieurs images.
\end{questions}
\editionRetenir{Un vecteur vitesse nul à un instant ne signifie pas une somme des forces nulle.}

\editionPage{defi-cor-14-1}{Solution du défi D14}{Une galerie à mesurer}{Une démarche possible ; d’autres rédactions sont recevables.}
\schema[corrige]{echo}
\begin{questions}
\item $d_1=71,4$ m ; $d_2=72,08$ m, soit environ 72,1 m.
\item L’écart est 0,68 m ; la lecture du retard, la position de la source ou la température peuvent contribuer.
\item Comparer les échos à plusieurs positions, repérer un trajet cohérent avec la géométrie et vérifier les réflexions sur les parois.
\end{questions}
\editionRetenir{Une distance sonore dépend du trajet identifié, pas seulement de la formule distance = vitesse fois durée.}

\editionPage{defi-cor-15-1}{Solution du défi D15}{Une note mal identifiée}{Une démarche possible ; d’autres rédactions sont recevables.}
\begin{questions}
\item $T=9,09/4=\qty{2,2725}{ms}$ ; $f\simeq\qty{440}{Hz}$.
\item Pour un signal sinusoïdal, maximum-minimum correspond à une demi-période. Prendre cette durée pour $T$ double $1/T$.
\item Les hauteurs sont identiques, les timbres peuvent différer. L’intensité exige une information sur l’amplitude et la chaîne de mesure.
\end{questions}
\editionRetenir{Avant de calculer une fréquence, vérifier le motif complet et l’unité du temps.}

\editionPage{defi-cor-16-1}{Solution du défi D16}{Identifier un éclairage}{Une démarche possible ; d’autres rédactions sont recevables.}
\begin{questions}
\item La première montre une émission de raies sur un fond ; la seconde est compatible avec un corps chaud dense dont la température augmente.
\item Non : au moins trois radiations sont présentes, en plus du fond.
\item Une table de spectres de référence et une échelle étalonnée permettraient de comparer plusieurs raies, avec la précision de mesure.
\end{questions}
\editionRetenir{La couleur apparente ne remplace pas l’analyse de la composition spectrale.}

\editionPage{defi-cor-17-1}{Solution du défi D17}{Deux couleurs dans le verre}{Une démarche possible ; d’autres rédactions sont recevables.}
\begin{questions}
\item $r_R=\arcsin(\sin40^\circ/1,50)\simeq25,4^\circ$ ; $r_V\simeq24,8^\circ$.
\item Le violet, car son angle avec la normale est plus petit.
\item Les directions de sortie diffèrent après les deux interfaces ; à distance, la séparation spatiale peut devenir observable.
\end{questions}
\editionRetenir{La dispersion s’explique par une dépendance de l’indice à la longueur d’onde.}

\editionPage{defi-cor-18-1}{Solution du défi D18}{Agrandir une projection}{Une démarche possible ; d’autres rédactions sont recevables.}
\schema[corrige]{lentille-agrandie}
\begin{questions}
\item La construction donne une image à 30 cm après la lentille, de taille 4,0 cm et renversée.
\item $\gamma=-4,0/2,0=-2,0$ ; par Thalès, $|\gamma|=30/15=2,0$, cohérent avec le renversement.
\item Lorsque l’objet était placé à 30 cm de la lentille, l’image mesurait 1,0 cm et se formait à 15 cm derrière la lentille. Rapprocher l’objet de F, tout en restant au-delà de F, produit ici une image agrandie mais demande davantage de recul pour l’écran.
\end{questions}
\editionRetenir{Pour projeter une image agrandie, il faut anticiper à la fois le placement de l’objet et le recul de l’écran.}

\editionPage{defi-cor-19-1}{Solution du défi D19}{Un circuit à diagnostiquer}{Une démarche possible ; d’autres rédactions sont recevables.}
\schema[corrige]{parallele-defi}
\begin{questions}
\item $I_1=6,0/200=30$ mA et $I_2=6,0/300=20$ mA.
\item $I=30+20=\qty{50}{mA}$ : l’indication est cohérente.
\item L’ampèremètre doit être en série dans la branche principale. Le placer seul en dérivation sur le générateur créerait un risque de court-circuit.
\end{questions}
\editionRetenir{Une mesure électrique demande à la fois une loi adaptée et un branchement correct.}

\editionPage{defi-cor-20-1}{Solution du défi D20}{Du code à l’alerte}{Une démarche possible ; d’autres rédactions sont recevables.}
\schema[corrige]{pont-capteur}
\begin{questions}
\item $U=5,00\times400/1000=2,00$ V ; $R_c=10,0\times2,00/(5,00-2,00)=6,67$ kΩ.
\item $\theta\simeq33,3\,^{\circ}\mathrm C$, donc l’alerte se déclenche dans ce modèle.
\item Rejeter les tensions proches de 5,00 V qui rendent le quotient instable ; vérifier que la température calculée reste dans le domaine 0 à 40 °C. Tester aussi les connexions et les valeurs limites.
\end{questions}
\editionRetenir{Une chaîne de mesure doit être validée à chaque étape : capteur, circuit, acquisition et calcul.}

